\documentclass[12pt,a4paper,oneside]{article}
\usepackage[brazil]{babel}
\usepackage[utf8]{inputenc}
\usepackage{olymp}


\newlength{\exmpwidouf}
\exmpwidinf=10cm
\exmpwidouf=6cm


\setcounter{problem}{3}

\begin{document}

\begin{problem}{Ordenação de capivaras}{ordempeso}{2}

Maria é bolsista de IC do curso de Zootecnia e precisa de sua ajuda. Durante sua pesquisa, coletou dados de várias capivaras da UFV e fez uma tabela com nome (que ela batizou), matrícula (que grampeou nas orelhas), e ainda peso e altura (que mediu após banhos na lagoa). Veja abaixo uma parte desta tabela.

\includegraphics[scale=0.75]{images/exercise_98_0.png}

Ela precisa ordenar as capivaras pelo peso.

% \Notes

% Você pode criar uma função para o cálculo do fatorial.

\InputFile

A entrada é iniciada por uma linha contendo o número $C$ de capivaras da tabela de Maria. Em seguida existem $C$ linhas, cada uma com 4 dados de uma capivara: $N$ $M$ $P$ $A$, respectivamente o nome da capivara, sua matrícula, seu peso e sua altura (em cm). Restrições: $N$ é um string contendo apenas letras minúsculas (no máximo 15) e os demais dados são números inteiros, com $1 \leq C \leq 100$.

\OutputFile

Seu programa deve escrever $C$ linhas na saída, listando os dados das capivaras por ordem de peso (não haverá empate). Os dados de cada capivara devem ser escritos no formato $N;M;P;A$, ou seja, separados por ponto-e-vírgula. Este é o formato CSV, que pode ser aberto por planilhas eletrônicas, como Excel, Calc e Numbers.

% \Notes
% Preste atenção aos tipos das variáveis, pois $F(90) \approx 2^{61}$.

\Example

\begin{example}
\exmp{
3
huguinho 110 80 30
zezinho 105 90 50 
luisinho 190 70 40
}{
luisinho;190;70;40
huguinho;110;80;30
zezinho;105;90;50 
}%
\end{example}

\begin{example}
\exmp{
7
chaves 8 75 160
florinda 14 70 165
girafales 999 98 250
barriga 1000 150 175
nhonho 1001 145 174
chiquinha 72 50 145
clotilde 71 74 170
}{
chiquinha;72;50;145
florinda;14;70;165
clotilde;71;74;170
chaves;8;75;160
girafales;999;98;250
nhonho;1001;145;174
barriga;1000;150;175
}%
\end{example}

% \textit{Problema baseado na questão ``A idade de dona Mônica'' da OBI2019 fase local.}

\end{problem}

\end{document}
